\uuid{7NZQ}
\exo7id{3473}
\titre{exo7 3473}
\auteur{quercia}
\organisation{exo7}
\datecreate{2010-03-10}
\isIndication{false}
\isCorrection{false}
\chapitre{Déterminant, système linéaire}
\sousChapitre{Système linéaire, rang}
\cpge{mpsi-mp2i.calcul-matriciel-systemes.a, pcsi.calcul-matriciel-systemes-lineaires.a, ptsi.calcul-matriciel-systemes-lineaires.a, mpsi-mp2i.matrices-applications-lineaires.b, pcsi.matrices-applications-lineaires.b, ptsi.matrices-applications-lineaires.b, mpsi-mp2i.nombres-complexes.a, pcsi.nombres-complexes.a, ptsi.nombres-complexes.a, bcpst1.nombres-complexes.a, bcpst1.matrices}
\module{Algèbre}
\niveau{L2}
\difficulte{}

\contenu{
\question{
Soit $M \in \mathcal{M}_n(\C)$ de rang 1. On écrit $M = P + iQ$ avec $P,Q \in \mathcal{M}_n(\R)$.
Montrer que rg$(P) \le 2$.
}
}
